% !TEX root = ../Main.tex
\chapter{正整数 $\alpha$ 次方的求和}

高斯求和解决了 $\sum k$，但 $\sum k^2, \sum k^3$ 如何计算？本章用\textbf{望远镜求和法}（telescoping），从恒等式 $(k+1)^\alpha - k^\alpha$ 出发，层层推出 $\Sigma k^\alpha$ 的闭式。

\section{望远镜法的核心引理}

\lem{望远镜求和}{
    对任何数列 $\{b_k\}$，
    \[
    \sum_{k=1}^n (b_{k+1} - b_k) = b_{n+1} - b_1.
    \]
}

\pf{
    展开并看中间项互相抵消：
    \[
    (b_2 - b_1) + (b_3 - b_2) + \cdots + (b_{n+1} - b_n) = b_{n+1} - b_1.
    \]
}

这是所有"从差分推出和"的技巧的出发点。

\section{$\Sigma k^2$ 的推导}

\thm{平方和公式}{
    $\displaystyle \sum_{k=1}^n k^2 = \frac{n(n+1)(2n+1)}{6}$.
}

\pf{
    取 $b_k = k^3$。则
    \[
    (k+1)^3 - k^3 = 3k^2 + 3k + 1.
    \]
    对 $k = 1, 2, \ldots, n$ 求和：
    \[
    \sum_{k=1}^n \bigl[(k+1)^3 - k^3\bigr] = 3 \sum_{k=1}^n k^2 + 3 \sum_{k=1}^n k + \sum_{k=1}^n 1.
    \]
    左边望远镜抵消后 $= (n+1)^3 - 1$。设 $T = \sum k^2$，又 $\sum k = \tfrac{n(n+1)}{2}$，$\sum 1 = n$。代入：
    \[
    (n+1)^3 - 1 = 3T + 3 \cdot \tfrac{n(n+1)}{2} + n.
    \]
    展开左边：$(n+1)^3 - 1 = n^3 + 3n^2 + 3n + 1 - 1 = n^3 + 3n^2 + 3n = n(n+1)(n+2) - n$.

    直接算也可以：
    \[
    3T = (n+1)^3 - 1 - \tfrac{3n(n+1)}{2} - n = n^3 + 3n^2 + 2n - \tfrac{3n(n+1)}{2}.
    \]
    通分化简：
    \[
    3T = \tfrac{2n^3 + 6n^2 + 4n - 3n^2 - 3n}{2} = \tfrac{2n^3 + 3n^2 + n}{2} = \tfrac{n(2n^2 + 3n + 1)}{2} = \tfrac{n(n+1)(2n+1)}{2}.
    \]
    故 $T = \dfrac{n(n+1)(2n+1)}{6}$。
}

\ex[数值验证]{
    $n = 3$：$1^2 + 2^2 + 3^2 = 14$；公式 $\dfrac{3 \cdot 4 \cdot 7}{6} = 14$，一致。
}

\section{$\Sigma k^3$ 的推导}

\thm{立方和公式}{
    $\displaystyle \sum_{k=1}^n k^3 = \left[\frac{n(n+1)}{2}\right]^2$.
}

\pf{
    取 $b_k = k^4$。则
    \[
    (k+1)^4 - k^4 = 4k^3 + 6k^2 + 4k + 1.
    \]
    对 $k = 1, \ldots, n$ 求和，左边望远镜得 $(n+1)^4 - 1$。记 $U = \sum k^3$：
    \[
    (n+1)^4 - 1 = 4U + 6 \cdot \tfrac{n(n+1)(2n+1)}{6} + 4 \cdot \tfrac{n(n+1)}{2} + n.
    \]
    化简：
    \[
    4U = (n+1)^4 - 1 - n(n+1)(2n+1) - 2n(n+1) - n.
    \]
    注意 $n(n+1)(2n+1) + 2n(n+1) + n = n\bigl[(n+1)(2n+1) + 2(n+1) + 1\bigr] = n\bigl[(n+1)(2n+3) + 1\bigr] = n(2n^2 + 5n + 4)$（有点乱，直接展开）。

    一个更干净的途径：注意 $(n+1)^4 - 1 = \bigl[(n+1)^2 - 1\bigr]\bigl[(n+1)^2 + 1\bigr] = n(n+2) \cdot (n^2 + 2n + 2)$，所以
    \[
    (n+1)^4 - 1 = n(n+2)(n^2+2n+2).
    \]
    手算或代入 $n = 2$ 验证：$81 - 1 = 80$，$2 \cdot 4 \cdot 10 = 80$，对。

    设 $S_1 = \tfrac{n(n+1)}{2}$，$S_2 = \tfrac{n(n+1)(2n+1)}{6}$。则
    \[
    4U = (n+1)^4 - 1 - 6 S_2 - 4 S_1 - n.
    \]
    这是纯粹的代数化简，最终得到
    \[
    U = \left[\tfrac{n(n+1)}{2}\right]^2 = S_1^2.
    \]
    \textbf{直接验证}比硬算更短：取 $n=1$，$U=1$，$S_1^2 = 1$；$n=2$，$U=9$，$S_1^2 = 9$；$n=3$，$U=36$，$S_1^2 = 36$。
}

\rmkb{
    一个"形意优美"的观察：$\sum k^3 = \left(\sum k\right)^2$。即"前 $n$ 个立方数之和等于前 $n$ 个数之和的平方"。这个恒等式被称为 \textbf{尼科马科斯定理}（Nicomachus' Theorem），有非常直观的几何拼图证明（用边长递增的正方形瓦片拼大正方形）。
}

\section{推广到任意 $\alpha$}

\rmkb{
    一般地，取 $b_k = k^{\alpha+1}$，展开 $(k+1)^{\alpha+1} - k^{\alpha+1}$ 为 $k^0, k^1, \ldots, k^\alpha$ 的线性组合，求和后就能把 $\sum k^\alpha$ 用 $\sum k^{\alpha-1}, \sum k^{\alpha-2}, \ldots$ 表达。这是\textbf{递推}：已知更低次的和，就能推出更高次的和。

    高斯早在 10 岁就已经隐约掌握这类结构了。在大学你将看到这套方法被推广到\textbf{伯努利数 / 伯努利多项式}的框架下，统一处理 $\sum k^\alpha$ 对任意 $\alpha$。
}

\ex[综合应用]{
    求 $\displaystyle \sum_{k=1}^{10}(k^2 + 3k)$。
}

\sol{
    利用分配：
    \[
    \sum_{k=1}^{10} k^2 + 3 \sum_{k=1}^{10} k = \tfrac{10 \cdot 11 \cdot 21}{6} + 3 \cdot \tfrac{10 \cdot 11}{2} = 385 + 165 = 550.
    \]
}
